\begin{figure}[htbp]
	\centering
	\begin{tikzpicture}[>=Stealth, line width=0.8pt, scale=1.0]
		
		\coordinate (O) at (0,0);
		\coordinate (f) at ({1.5*cos(30)},{1.5*sin(30)});
		% C: mu_z=1.5, |f|=1.5 -> |OC|=1.0
		% C = (mu_z/|f|^2)*(fy,-fx) = (0.5, -0.866)
		\coordinate (C) at (0.500, -0.866);
		
		% Retta perpendicolare a f tratteggiata con freccia verso infinito
		\draw[dashed, gray!60]
		({-0.3*cos(-60)},{-0.3*sin(-60)}) --
		({2.2*cos(-60)}, {2.2*sin(-60)});
		\draw[->, dashed, gray!60]
		({2.2*cos(-60)},{2.2*sin(-60)}) --
		({2.7*cos(-60)},{2.7*sin(-60)})
		node[right=1pt] {\small$\fb\to \mathbf{0}$};
		
		% --- Sistema in O: forza + coppia ---
		% Vettore f in O
		\draw[->, thick] (O) -- (f) node[above right=1pt] {$\fb$};
		% Coppia mu_z in O (freccia semicircolare antioraria)
		\draw[<-, thick] (-0.55, 0) arc[start angle=180, end angle=60, radius=0.55];
		\node[left=2pt] at (-0.55, 0.3) {$\mu_z(O)$};
		
		% Angolo retto
		\draw[thin]
		({0.15*cos(30)},{0.15*sin(30)}) --
		({0.15*cos(30)+0.15*cos(-60)},{0.15*sin(30)+0.15*sin(-60)}) --
		({0.15*cos(-60)},{0.15*sin(-60)});
		
		% Quota |OC| = mu_z/|f| (traslata a sinistra oltre la tratteggiata)
		\draw[<->, thin]
		(-0.433, -0.250) --
		(0.067, -1.116)
		node[midway, left=3pt] {\small$\dfrac{|\mu_z(O)|}{|\fb|}$};
		
		% Punto O e punto C
		\fill (O) circle (2pt) node[below left=-4pt] {$O$};
		\fill[white] (C) circle (3pt);
		\draw[thick] (C) circle (3pt);
		\node[below right=2pt] at (C) {$C$};
		
		% --- Sistema equivalente in C: solo forza ---
		% Vettore f in C (traslato)
		\draw[->, thick] (C) -- ++({1.5*cos(30)},{1.5*sin(30)})
		node[above right=1pt] {$\fb$};
		% Nessuna coppia (mu_z(C)=0)
		%				\node[below left=2pt] at ({0.5+1.5*cos(30)/2},{-0.866+1.5*sin(30)/2})
		%				{\small$\mu_z(C){=}0$};
		
		% Freccia tratteggiata che collega i due sistemi
		\draw[dashed, gray!40] (O) -- (C);
		
	\end{tikzpicture}
	\caption{Centro di riduzione\index{centro di riduzione|seealso{asse di sollecitazione}} nel caso piano. Il sistema equivalente
		a una forza $\fb$ applicata in $O$ con momento risultante $\mu_z(O)$
		si riduce alla sola forza $\fb$ applicata nel punto $C$ della
		perpendicolare a $\fb$ passante per $O$, a distanza
		$|\protect\overrightarrow{OC}| = |\mu_z(O)|/|\fb|$.
		Al tendere di $|\fb|\to 0$ a parità di $\mu_z(O)$ il centro
		si allontana all'infinito: una coppia pura non ammette centro
		di riduzione.}
	\label{fig:asse_sollecitazione}
\end{figure}
